\sscp{\tsc{Ambon-Seram} and \tsc{Seram-Tanimbar-Bomberai}}{07-03-03}
\tsc{Ambon-Seram}, \tsc{Seram-Tanimbar-Bomberai}, and Banda have
a merger of *d/*z in most words.
The outcome in Proto-Seram-Tanimbar-Bomberai was
probably **d, as Yamdena has *d > \it{r,
d} and *z > \it{d}.
However, Banda has *d/*z > \it{r},
and the outcome in \tsc{Ambon-Seram} is more diverse: \tsc{Nuclear
Ambon-Seram} has *d/*z > **l (reflexes as \it{r} or \it{d} are due to later changes; \srf{10-01}),
%\footnote{
		%Reflexes of *d/*z in \tsc{Nuclear
		%Ambon-Seram} as \it{r} or \it{d} are due to later changes
		%of **l (\srf{10-01}).}
\tsc{Patakai-Manusela} has *d/*z > **h,
and \tsc{Atamanu} has *d/*z > **r.
All these \tsc{Ambon-Seram} reflexes can probably be explained
by positing *d/*z > **r for Proto-Ambon-Seram.

While it is not unlikely that **r in \tsc{Ambon-Seram}
and Banda went through intermediate **d,
it is not certain.
The merger of *d/*z in these branches could have
had a different quality,
such as [z] and/or [ʤ],
in which case the merger of *d/*z in \tsc{Ambon-Seram},
Banda, and \tsc{Seram-Tanimbar-Bomberai} would not be a shared innovation.
In any case, if merger of *d/*z occurred once in a common ancestor
of \tsc{Ambon-Seram}, Banda and \tsc{Seram-Tanimbar-Bomberai}, the evidence indicates
that Proto-Ambon-Seram subsequently underwent a change independent of \tsc{Seram-Tanimbar-Bomberai}.
See \srf{13-04} for more discussion.

Examples of the merger are given in \trf{07-07}.
Note that many \tsc{Nuclear Ambon-Seram} languages have subsequent
changes to *d/*z > (**r >) **l conditioned by the adjacent vowels (see \srf{10-01}).
Some members of \tsc{Tanimbar-Bomberai} have conditioned reflexes
of *d when the following syllable contains /r/.
That is, in the environment /{\gap}Vr, e.g.
*d in *duRi ‘thorn,
bone’ has different reflexes to *d in *duha ‘two’(see \srf{11-05-01}).
Most languages of \tsc{Sula-Buru} have *z/*y > **y,
but different reflexes of *d,
meaning we cannot posit merger of *d/*z for \tsc{Sula-Buru}.

%\begin{table}[p]\centering\tcp{\tsc{Ambon-Seram} and Banda *d and *z}{07-07a}
\begin{table}[p]\centering
\caption[Central Maluku *d and *z]{Central Maluku *d and *z\su{†}}\label{tab:07-07}
\begin{adjustbox}{totalheight=\textheight}
	\begin{threeparttable} \st{4pt}
	\begin{tabular}{llll|lll}\lsptoprule
*gloss	&	path	&	far	&	rain	&	leaf	&	two	&	thorn, bone\\
PMP	&	*\tbr{z}alan\su{‡}	&	*\tbr{z}auq	&	*qu\tbr{z}an	&	*\tbr{d}ahun	&	*\tbr{d}uha	&	*\tbr{d}uRi\\ \midrule
%\mc{7}{c}{\tsc{Ambon-Seram}}\\	\midrule
Kayeli	&	{\hr}\tbr{l}alan	&	{\hr}hat\tbr{l}au	&	{\hr}u\tbr{l}an	&	\cg{(au tetun)}	&	{\hr}\tbr{l}ua	&	{\hr}\tbr{l}oli-n\\
Boano	&	{\hr}\tbr{l}alan-e	&	{\hr}-\tbr{l}ou	&	{\hr}u\tbr{l}an-e	&		&	{\hr}\tbr{l}ua	&	{\hr}\tbr{l}uli\\
Luhu	&	{\hr}\tbr{l}alan	&	{\hr}-\tbr{l}au	&	{\hr}u\tbr{l}an-e	&	\cg{(wata)}	&	{\hr}\tbr{l}ua	&	{\hr}\tbr{l}uli\\
Larike	&	{\hr}\tbr{l}alan-u	&	{\hr}\tbr{l}ou	&	{\hr}u\tbr{\und{d}}an-u	&	{\hr}\tbr{l}ou	&	{\hr}\tbr{\und{d}}ua	&	{\hr}\tbr{\und{d}}udi\\
Asilulu	&	{\hr}\tbr{l}alan	&	{\hr}\tbr{l}au	&	{\hr}u\tbr{l}a, u\tbr{l}an	&	{\hr}\tbr{l}aun	&	{\hr}\tbr{l}ua	&	{\hr}\tbr{l}uli-\\
Paulohi	&	{\hr}\tbr{l}aa kine	&	{\hr}e\tbr{l}au	&	\cg{(kia)}	&	{\hr}\tbr{l}au	&	{\hr}\tbr{\und{ʀ}}ua	&	{\hr}\tbr{\und{ʀ}}iʀi-a\\
Kaibobo	&	{\hr}\tbr{l}alan	&	{\hr}\tbr{l}au	&	{\hr}u\tbr{\und{r}}an	&	{\hr}\tbr{l}au	&	{\hr}\tbr{\und{r}}ua	&	{\hr}\tbr{\und{r}}uri\\
Kamarian	&	{\hr}\tbr{l}alan	&	{\hr}\tbr{l}au	&	{\hr}u\tbr{\und{r}}an	&	{\hr}\tbr{l}au	&	{\hr}\tbr{\und{r}}ua	&	{\hr}\tbr{\und{r}}uri\\
Haruku	&	{\hr}\tbr{l}alan	&	{\hr}\tbr{l}au	&	{\hr}u\tbr{\und{r}}an	&	{\hr}\tbr{l}au	&	{\hr}\tbr{\und{r}}ua	&	{\hr}\tbr{\und{r}}uri\\
Saparua	&	{\hr}\tbr{l}alan-o	&	{\hr}\tbr{l}au	&	\cg{(kial)}	&	{\hr}\tbr{l}au	&	{\hr}\tbr{l}ua	&	{\hr}\tbr{l}ilin\\
Nusalaut	&	{\hr}\tbr{l}alan-o	&	{\hr}\tbr{l}au	&	\cg{(kial)}	&	{\hr}\tbr{l}au	&	{\hr}\tbr{\und{r}}ua	&	{\hr}\tbr{\und{r}}iriɲ-o\\
Sepa	&	{\hr}\tbr{l}aa tina	&	{\hr}\tbr{l}au	&	\cg{(kiya)}	&	{\hr}\tbr{l}aun	&	{\hr}\tbr{\und{r}}ua	&	\cg{(toi)}\\
N. Wemale	&	{\hr}\tbr{l}arina	&	{\hr}e\tbr{l}aute	&	{\hr}u\tbr{l}an-e	&	{\hr}\tbr{l}auni	&	{\hr}\tbr{l}ua	&	{\hr}\tbr{l}uli\\
N. Alune	&	{\hr}\tbr{l}alan-e	&	{\hr}\tbr{l}au, \tbr{l}aukwe	&	{\hr}u\tbr{l}an-e	&	{\hr}\tbr{l}oini	&	{\hr}\tbr{l}ua	&	{\hr}\tbr{l}uli, \tbr{l}uri\\
Yalahatan	&	{\hr}\tbr{r}aran-e	&	{\hr}a\tbr{r}aʔu-e	&	{\hr}u\tbr{r}a, u\tbr{r}an-e	&	{\hr}\tbr{l}ain	&	{\hr}\tbr{r}ua	&	{\hr}\tbr{r}uri-n\\
S. Nuaulu	&	\cg{(arena)}	&	\cg{(hai-nau)} \su{‡}	&	{\hr}uan-e	&	\cg{(totu-e)}	&	{\hr}ua	&	{\hr}uni-e\\
Manusela	&	{\hr}\tbr{h}ala	&	\cg{(meleke)}	&	\cg{(roa)}	&	\cg{(ai totu)}	&	{\hr}\tbr{h}ua	&	{\hr}\tbr{h}uli-kata\\
Liana-Seti	&	{\hr}\tbr{l}ala	&	\cg{(koila)}	&	\cg{(roa)}	&	{\hr}\tbr{l}ana	&	{\hr}en-\tbr{l}ua	&	{\hr}en-\tbr{l}uil-a\\ \midrule
Banda	&	{\hr}e\tbr{r}an	&	{\hr}\tbr{r}au	&	{\hr}u\tbr{r}en	&	{\hr}\tbr{r}ano	&	{\hr}\tbr{r}uo	&	{\hr}\tbr{r}iri\\
%\mc{7}{c}{\tsc{Seram-Tanimbar-Bomberai}}\\	\midrule
Bobot	&	{\hr}\tbr{l}olan	&		&	{\hr}u\tbr{l}an	&	{\hr}ai \tbr{l}an	&	{\hr}\tbr{l}ua	&	{\hr}-\tbr{l}ui-\\
Masiwang	&	{\hr}\tbr{l}olon	&		&	\cg{(roti)}	&	{\hr}han \tbr{l}an	&	{\hr}ain \tbr{l}ua	&	{\hr}\tbr{l}usi\\
Geser-Gorom	&	{\hr}\tbr{l}alan	&	{\hr}\tbr{r}au	&	{\hr}u\tbr{r}an	&	{\hr}\tbr{r}u	&	{\hr}\tbr{r}oti	&	{\hr}\tbr{r}uri\\
Watubela	&		&		&	\cg{(kadama)}	&	{\hr}\tbr{l}ueʔ	&	{\hr}\tbr{l}ua	&	{\hr}\tbr{l}uli\\
Kowiai	&	{\hr}\tbr{l}araʔai	&	{\hr}\tbr{s}or	&	\cg{(oma)}	&	{\hr}\tbr{r}owa	&	{\hr}\tbr{r}uwet	&	{\hr}\tbr{r}ur\\
Teor	&		&		&	‹hu\tbr{r}ani›	&	\cg{(‹chafen›)}	&	‹\tbr{r}úa›	&	\cg{(‹urut›)}\\
Kur	&		&	\cg{(anloin)}	&	‹oe\tbr{r}an›	&		&	{\hr}\tbr{r}ua	&	\\
Yamdena	&	{\hr}\tbr{d}alnen-e\su{Δ}	&	{\hr}\tbr{d}o-do	&	{\hr}u\tbr{d}an	&	{\hr}\tbr{d}on-e	&	{\hr}\tbr{d}u	&	{\hr}\tbr{d}uri\\
Fordata	&	\cg{(liŋaʔan)}	&	{\hr}\tbr{r}aroa	&	\cg{(daʔut)}	&	{\hr}\tbr{r}oa	&	{\hr}\tbr{r}ua	&	{\hr}\tbr{l}uri-n\\
Kei	&	\cg{(laŋay)}	&	{\hr}\tbr{r}o-ro	&	\cg{(doʔot)}	&	{\hr}ron	&	{\hr}\tbr{r}u	&	{\hr}\tbr{l}uri-n\\
Uruangnirin	&	{\hr}\tbr{n}enmatan	&	\cg{(olat)}	&	{\hr}u\tbr{n}an	&	{\hr}ai \tbr{n}on-in	&	{\hr}\tbr{n}ua	&	{\hr}\tbr{n}uri-n\\
Onin	&	{\hr}\tbr{n}aman	&	\cg{(bair)}	&	{\hr}maro u\tbr{n}an	&	{\hr}wi \tbr{n}on	&	{\hr}\tbr{n}ua	&	{\hr}\tbr{r}uri-r\\
Sekar	&	{\hr}\tbr{n}anam	&	\cg{(bair)}	&	\cg{(yagin)}	&	\cg{(kafak-in)}	&	{\hr}\tbr{n}ua	&	{\hr}\tbr{r}uri-r\\
Arguni	&	{\hr}\tbr{r}árin	&	{\hr}\tbr{d}ówan	&	\cg{(úmun)}	&	\cg{(náne)}	&	{\hr}\tbr{r}u	&	\cg{(tór)}\\
Bedoanas	&	{\hr}\tbr{r}árin	&	{\hr}\tbr{nd}ewen\su{‖}	&	{\hr}á\tbr{r}i	&	{\hr}\tbr{r}ína	&	{\hr}\tbr{r}uyu	&	\cg{(tóʔor)}\\
Erokwanas	&	{\hr}\tbr{r}órin	&	{\hr}\tbr{d}áwen	&	\cg{(úmin)}	&	\cg{(manda-}	&	{\hr}\tbr{r}u	&	\cg{(toʔar)}\\
					&										&										&							&	\cg{-nánia)}	&							&							\\
Irarutu	&	{\hr}\tbr{r}arn	&	{\hr}ne\tbr{r}o	&	\cg{(syem)}	&	{\hr}\tbr{r}o	&	{\hr}\tbr{r}ifo	&	{\hr}\tbr{r}ur\\
Kuri	&	{\hr}\tbr{r}arn	&	{\hr}ge\tbr{r}o	&	\cg{(naun)}	&	{\hr}rue\tbr{r}o	&	{\hr}\tbr{r}u	&	{\hr}\tbr{r}ur\\
		\lspbottomrule
	\end{tabular}
		\begin{tablenotes}\tnsize
			\item[†]	{Reflexes in \tsc{Ambon-Seram} languages which are conditioned
								by an adjacent high vowel (\srf{10-01}) are underlined.}
			\item[‡]	{South Nuaulu \it{hai-nau} ‘far’ is from *lahud ‘seaward, sea’, not *zauq.}
			\item[Δ]	{\citet{MettlerMettler1996} list multiple words for ‘road,
								path’: \it{dalamten-e, dalamten-ar, dalnen-e, dalnyen-e,} and \it{dalŋen}.
								Sekar, Onin, and some varieties of Irarutu also show final *n > \it{m}.
								Onin \it{naman} has consonant metathesis from earlier **nanam.}
			\item[‖]	{Bedoanas [nd] appears to be a variant of /d/,
								at least word initially.
								This is also attested in *dalij ‘buttress root’ > \it{dara} {\tl} \it{ndarao} ‘root’.
								(Arguni \it{dare}, Erokwanas \it{dara} ‘root’).}
		\end{tablenotes}
	\end{threeparttable}
\end{adjustbox}
\end{table}

\tsc{Ambon-Seram}, Banda, and nearly all \tsc{Seram-Tanimbar-Bomberai}
languages also usually have lenition of *p > **f/**ɸ,
often with subsequent further lenition to \it{h, {\0}}.
However, Kur retains *p unchanged in all positions (\srf{11-04}),
and there are sporadic retentions of *p = \it{p} in \tsc{Ambon-Seram}.
Examples of *p = \it{p} in \tsc{Ambon-Seram} include: *\tbr{p}əñu > Alune
\it{\tbr{p}one} ‘(land) tortoise’, (cf.
Alune \it{benu} ‘sea turtle’),
Larike \it{\tbr{p}one{\tl}\tbr{p}one} ‘(land) tortoise’,
as well as *\tbr{p}a- \tsc{causative} > South Nuaulu \it{a\tbr{p}a-,
a\tbr{p}u}; Asilulu \it{\tbr{p}a-, a-},
and Larike \it{\tbr{p}a-, a-}.

Yamdena and Banda also retain *p unchanged word finally.
Yamdena examples include: *mali\tbr{p} > \it{mali\tbr{p}} ‘laugh’
and *ma-qudi\tbr{p} > \it{n-mori\tbr{p}} ‘live’. The single
Banda example in the available data is *qatə\tbr{p} > \it{ato\tbr{p}} ‘thatch, roof’.
Banda has initial and medial *p > {\0}.
This could have passed through intermediate **f/**ɸ in which
case it would be a change shared with \tsc{Ambon-Seram}
and most of \tsc{Seram-Tanimbar-Bomberai},
but it could have had an alternate intermediate stage,
such as [ʔ], in which case the reflexes of *p would not be shared
between Banda and other languages of the region.

Weakening of *p is an areal change which also occurs
in \tsc{Aru}, \tsc{Timor-Babar}, and SHWNG.
Within \tsc{Sula-Buru} only Ambelau has *p > \it{f},
with other \tsc{Sula-Buru} languages retaining *p = \it{p}.

